\section{Metodología de Implementación y Ruta Crítica}

La evaluación de implementación determinó que el proyecto se ejecutaría bajo un marco escalonado y flexible, priorizando la entrega progresiva de valor a lo largo de una duración de 6 meses. Esto permitió validar resultados de manera temprana en la tienda de la calle Deán Valdivia.

\subsection{Metodología Ágil (Scrum)}
Se aplicó estrictamente la metodología ágil Scrum, considerada ideal para entornos digitales por su capacidad de entregar resultados tangibles en ciclos cortos (Sprints). Este enfoque permitió que Textiles y Confecciones Atlas viera mejoras reales en su control de inventarios desde el segundo mes. 
Los roles definidos fueron:
\begin{itemize}
    \item \textbf{Product Owner (Gerencia General):} Encargado de priorizar las funciones del sistema, ordenando por ejemplo la carga prioritaria de los productos más vendidos antes que los de menor rotación.
    \item \textbf{Development Team:} Conformado por el proveedor tecnológico (encargados de configurar el software POS) y el equipo responsable del diseño de las campañas de marketing.
    \item \textbf{Product Backlog:} La lista priorizada que incluyó el registro digital de tallas, modelos, precios, catálogos y la base de datos de clientes institucionales.
\end{itemize}

\subsection{La Ruta Crítica}
La ruta crítica delineó las tareas esenciales ineludibles para cumplir con el lanzamiento (Go-Live) en el mes 6, proyectado para el cierre de julio de 2026:
\begin{enumerate}
    \item \textbf{Levantamiento de información técnica:} Realización del inventariado manual final indispensable para la migración de datos.
    \item \textbf{Selección de la plataforma POS/Ventas:} Elección de un sistema nativo en la nube compatible con los recursos actuales (la PC de la tienda).
    \item \textbf{Carga de datos de inventario:} Digitalización masiva de existencias de artículos deportivos (camisetas, balones, trofeos).
    \item \textbf{Configuración del módulo de Ventas (POS):} Integración fluida con el flujo diario de atención al cliente en mostrador.
    \item \textbf{Pruebas de sincronización:} Validación en entorno cerrado para asegurar que el stock descienda correctamente tras la simulación de cada venta.
    \item \textbf{Capacitación del personal:} Entrenamiento táctico en el uso del sistema y atención digital para vencer la barrera generacional.
    \item \textbf{Lanzamiento piloto controlado y oficial.}
\end{enumerate}

\subsection{Presupuesto y Análisis Costo-Beneficio}
El desarrollo y despliegue del proyecto fue apalancado por una inversión estratégica estructurada bajo el contrato de ProInnóvate. El Cuadro de Presupuesto Aprobado detalla un monto total de \textbf{S/ 57,143.00} distribuido de la siguiente forma:
\begin{itemize}
    \item \textbf{Aporte ProInnóvate (RNR):} S/ 40,000.00
    \item \textbf{Aporte Monetario (Entidad Ejecutora):} S/ 10,286.00
    \item \textbf{Aporte No Monetario (Entidad Ejecutora):} S/ 6,857.00
\end{itemize}

Los rubros principales incluyeron \textbf{Consultoría (S/ 49,886.00)} para la Implementación del Sistema de Gestión Comercial (configuración de inventarios y despliegue del Catálogo Digital) y el Pack Trimestral de Marketing Digital y Gestión de Redes (que incluyó estrategia de lanzamiento, redacción de publicaciones en Facebook/IG/TikTok y administración de Ads). Además, se destinaron fondos a Gastos de Difusión (Hito 3) y Honorarios del Coordinador General (S/ 6,857.00).

Este costo-beneficio validó que la inversión supera ampliamente el gasto inicial: los beneficios tangibles (reducción del 15\% de errores logísticos y tiempos de atención) y los beneficios intangibles (mejora drástica en la imagen de marca y alcance regional) mitigan el riesgo inminente de estancamiento operativo que habría llevado a la pérdida progresiva de clientes frente a la competencia digitalizada.
